%-------------------------
% Resume in Latex
% Based off of: https://github.com/jakegut/resume (Jake Gutierrez, MIT)
%------------------------

\documentclass[letterpaper,11pt]{article}

\usepackage{latexsym}
\usepackage[empty]{fullpage}
\usepackage{titlesec}
\usepackage{marvosym}
\usepackage[usenames,dvipsnames]{color}
\usepackage{verbatim}
\usepackage{enumitem}
\usepackage[hidelinks]{hyperref}
\usepackage{fancyhdr}
\usepackage[italian]{babel}
\usepackage{tabularx}
\usepackage[T1]{fontenc}
\usepackage[utf8]{inputenc}
\input{glyphtounicode}

\pagestyle{fancy}
\fancyhf{}
\fancyfoot{}
\renewcommand{\headrulewidth}{0pt}
\renewcommand{\footrulewidth}{0pt}

% A4 page, adjusted margins
\usepackage[a4paper,top=0.45in,bottom=0.45in,left=0.5in,right=0.5in]{geometry}

\urlstyle{same}

\raggedbottom
\raggedright
\setlength{\tabcolsep}{0in}

% Sections formatting
\titleformat{\section}{
  \vspace{-4pt}\scshape\raggedright\large
}{}{0em}{}[\color{black}\titlerule \vspace{-5pt}]

% Ensure that generate pdf is machine readable/ATS parsable
\pdfgentounicode=1

%-------------------------
% Custom commands
\newcommand{\resumeItem}[1]{
  \item\small{
    {#1 \vspace{-2pt}}
  }
}

\newcommand{\resumeSubheading}[4]{
  \vspace{-2pt}\item
    \begin{tabular*}{0.97\textwidth}[t]{l@{\extracolsep{\fill}}r}
      \textbf{#1} & #2 \\
      \textit{\small#3} & \textit{\small #4} \\
    \end{tabular*}\vspace{-7pt}
}

\newcommand{\resumeProjectHeading}[2]{
    \item
    \begin{tabular*}{0.97\textwidth}{l@{\extracolsep{\fill}}r}
      \small#1 & #2 \\
    \end{tabular*}\vspace{-7pt}
}

\renewcommand\labelitemii{$\vcenter{\hbox{\tiny$\bullet$}}$}

\newcommand{\resumeSubHeadingListStart}{\begin{itemize}[leftmargin=0.15in, label={}]}
\newcommand{\resumeSubHeadingListEnd}{\end{itemize}}
\newcommand{\resumeItemListStart}{\begin{itemize}}
\newcommand{\resumeItemListEnd}{\end{itemize}\vspace{-5pt}}

%-------------------------------------------
%%%%%%  RESUME STARTS HERE  %%%%%%%%%%%%%%%%%%%%%%%%%%%%


\begin{document}

%----------HEADING----------
\begin{center}
    \textbf{\Huge \scshape Lorenzo Baladier} \\ \vspace{2pt}
    \small \href{mailto:contact@lorenzo-baladier.fr}{\underline{contact@lorenzo-baladier.fr}} $|$
    \href{https://linkedin.com/in/lorenzo-baladier}{\underline{linkedin.com/in/lorenzo-baladier}} $|$
    \href{https://github.com/Loren844}{\underline{github.com/Loren844}} $|$
    \href{https://lorenzo-baladier.fr}{\underline{lorenzo-baladier.fr}} \\ \vspace{2pt}
    \textbf{\small Ingegnere software embedded $\cdot$  Cerco un posto a tempo indeterminato a Parigi da settembre 2027}
\end{center}

%-----------FORMAZIONE-----------
\section{Formazione}
  \resumeSubHeadingListStart
    \resumeSubheading
      {Polytech Nice Sophia -- Université Côte d'Azur}{Sophia Antipolis, FR}
      {Laurea magistrale in Ingegneria Informatica, IoT \& Sistemi Cyber-Fisici, apprendistato}{2024 -- 2027}
    \resumeSubheading
      {IUT Nice Côte d'Azur}{Sophia Antipolis, FR}
      {Diploma universitario in Informatica (DUT), apprendistato}{2022 -- 2024}
    \resumeSubheading
      {Lycée Pierre et Marie Curie}{Menton, FR}
      {Maturità scientifica francese, sezione EsaBac (doppio diploma franco-italiano)}{2019 -- 2022}
  \resumeSubHeadingListEnd

%-----------ESPERIENZA-----------
\section{Esperienza}
  \resumeSubHeadingListStart

    \resumeSubheading
      {Ingegnere software embedded \& SDV -- Apprendista}{Set. 2024 -- Oggi}
      {Ampere Software Technology}{Biot, FR}
      \resumeItemListStart
        \resumeItem{Progettazione e benchmark di una piattaforma di prototipazione basata sull'ambiente \textbf{Software-Defined Vehicle} di Google : \textbf{CarOS}}
        \resumeItem{Bring-up della piattaforma su \textbf{Raspberry Pi} e su dispositivi virtuali \textbf{Cuttlefish}; realizzazione delle pipeline \textbf{CI/CD}}
        \resumeItem{Sviluppo di casi d'uso end-to-end: comunicazione orientata ai servizi, diagnostica del veicolo ed elaborazione dati}
        \resumeItem{Riscrittura di un prototipo Python in uno strumento \textbf{Rust} modulare (CLI/TUI) che genera comandi di test \textbf{SOME/IP}, \textbf{Protobuf} a partire da GraphQL, testato con mock basati su trait}
      \resumeItemListEnd

    \resumeSubheading
      {Tirocinante in robotica \& visual servoing}{Mag. 2026 -- Ago. 2026}
      {DNIIT -- Da Nang Institute of Innovative Technology}{Da Nang, Vietnam}
      \resumeItemListStart
        \resumeItem{Visual servoing ad anello chiuso di un braccio \textbf{SCARA} industriale per il pick-and-place di oggetti in movimento}
        \resumeItem{Integrazione del rilevamento oggetti in tempo reale \textbf{YOLO} e della stima della posa con \textbf{ROS 2} e \textbf{MoveIt 2}}
        \resumeItem{Validazione della cinematica e del controllo PID in \textbf{MATLAB} e \textbf{Gazebo} prima del deployment sul robot reale -- \href{https://github.com/Loren844/project-10-dniit}{\underline{codice sorgente}}}
      \resumeItemListEnd

    \resumeSubheading
      {Sviluppatore software -- Apprendista}{Set. 2022 -- Lug. 2024}
      {ESI -- European Systems Integration}{Le Cannet, FR}
      \resumeItemListStart
        \resumeItem{Sviluppo di un'applicazione \textbf{Java} per l'esportazione in JSON di dati complessi da un database remoto}
        \resumeItem{Automazione dei test end-to-end con Selenium/TestNG, poi migrazione a \textbf{TypeScript/Cypress}}
      \resumeItemListEnd

  \resumeSubHeadingListEnd

%-----------PROGETTI-----------
\section{Progetti}
    \resumeSubHeadingListStart
      \resumeProjectHeading
          {\textbf{Emulatore CHIP-8} $|$ \emph{Rust, SDL2} -- \href{https://github.com/Loren844/rust-chip-8}{\underline{github.com/Loren844/rust-chip-8}}}{2025}
          \resumeItemListStart
            \resumeItem{Emulazione della macchina virtuale CHIP-8: ciclo fetch-decode-execute, memoria, registri, timer, display e tastiera}
          \resumeItemListEnd
      \resumeProjectHeading
          {\textbf{Go-Biking} $|$ \emph{C\#, SOAP/REST, ActiveMQ}}{Progetto universitario}
          \resumeItemListStart
            \resumeItem{Pianificatore di percorsi orientato ai servizi che combina API di bike sharing e di routing, con proxy di cache e messaggistica asincrona}
          \resumeItemListEnd
    \resumeSubHeadingListEnd

%-----------COMPETENZE-----------
\section{Competenze}
 \begin{itemize}[leftmargin=0.15in, label={}]
    \small{\item{
     \textbf{Linguaggi}{: Rust, C, C++, Python, Java, TypeScript, Bash} \\
     \textbf{Embedded \& automotive}{: Linux, Raspberry Pi, Cuttlefish, Software-Defined Vehicle, SOA, SOME/IP, Protobuf} \\
     \textbf{Strumenti}{: Git, GitLab CI/CD, Cypress, Selenium} \\
     \textbf{Lingue}{: Francese (madrelingua), Italiano (fluente, EsaBac), Inglese (TOEIC 955/990)}
    }}
 \end{itemize}

\end{document}
